\documentclass[12pt]{article}
\usepackage[utf8]{inputenc}
\usepackage[T1]{fontenc}
\usepackage{hyperref}
\title{Lower Burdekin Language Resources for K-12 Curriculum Development}
\author{Compiled by Claude Sonnet 5.5 (claude-sonnet-5-5) for Andr\'as Kornai}
\date{}
\begin{document}
\maketitle

\section{Introduction}
Lower Burdekin (ISO 639-3 xbb) is a geographical grouping of three extinct Aboriginal languages once spoken around the mouth of the Burdekin River in north Queensland. Only short wordlists survive, collected in the 19th century and published in volume 2 of \emph{The Australian Race} (1886). Almost nothing is known of the grammar. Gavan Breen analysed two of the lists and concluded they were different languages, neither Maric. He presumes one of them was Bindal (xbd). Others, such as Bowern, have suggested Bindal might be Maric. A curriculum using the language itself is not possible beyond a handful of recorded words.

\section{Primary Educational Resources}
None found.

\section{Linguistic and Reference Materials}
The Wikipedia article on the Lower Burdekin languages summarises the sources and the classification debate: \url{https://en.wikipedia.org/wiki/Lower_Burdekin_languages}. Glottolog lists a reference, Cunningham and Gorton 1886, which I saw only as a bibliographic record: \url{https://glottolog.org/resource/reference/id/106361}. The AIATSIS Austlang record E62 appeared in results; the search summary described it as covering a neighbouring language from Bowen north to the Burdekin River, and I did not verify its content: \url{https://collection.aiatsis.gov.au/austlang/language/e62}.

\section{Religious and Community Resources}
None found specific to xbb. The related Bindal language is covered under xbd.

\section{Digital and Audio Resources}
None found.

\section{Curriculum Development Context}
Any heritage programme would need to work with the Traditional Owner groups of the Burdekin and Bowen region (for example Bindal and Juru, as named in the search results), and with AIATSIS and archival wordlists. Material for neighbouring, better-documented languages could supplement a cultural-awareness curriculum.
\end{document}